\subsection{方法、候选资料与本项目新增计算的来源}
\label{app:source_provenance}
本文使用的自适应大邻域框架参考Ropke和Pisinger\cite{ropke2006alns}，任务列松弛和定价的一般方法参见文献\cite{barnhart1998branchprice}；具体物理参数与计算规则以题目附件为准。物理等价箱的身份匹配、飞机--电池解码、连续地形通信事件和固定日程资源核算的实现范围已在对应章节说明，不能把标准方法名称本身作为新算法贡献。

补充候选归档\cite{candidateArchive}提供过紧凑的箱组、机型、访问顺序以及中继坐标和高度。项目记录的原始文件名为\texttt{26研赛D代码结果最终版(1).zip}，SHA-256为
\begin{center}\begin{minipage}{.95\linewidth}\footnotesize\path{dffa94db1788b97661b40bd3359b5abcfdf0e017827d9d81720ba9f6989fb97c}\end{minipage}\end{center}
接收的候选先作为初始化或对照，随后在本项目的AEQD/DEM、$g=9.81$及当前中继转场模型下重新计算物理属性、重排全部执行时刻并独立核验。外部保存的目标值未被直接当作本项目的验收结果；但重新计算与核验也不等于原始候选构造由本文独立提出。

\begin{table}[htbp]\centering\small
\caption{候选来源与新增计算的职责分离}\label{tab:source_responsibility}
\begin{tabularx}{\linewidth}{p{3cm}XX}\toprule
对象 & 继承或引用的内容 & 本项目重新计算或新增的内容\\\midrule
问题二正式候选 & 已保存23架次基线；选定记录标记种子9201 & 独立物理重放、原始箱号匹配、食品/卫生箱交换的节能核对，以及后来明确预算的扰动与搜索试验\\
问题三初始化 & 补充归档中的部分运输模式与中继位置/高度 & 当前物理模型下重建、同模型中继局部重定位、整数联合排程、固定结构时间精修及完整通信验收\\
问题三约88秒改善 & 同一物理模型的改进前23项运输模式 & 保持运输模式不变，调整中继和资源同步关系；不将旧高度域变化混入纯搜索收益\\
问题四配置与稳定性 & 接受的上游联合日程及真实提供方记录 & 构建不可拆关系、完整枚举分区、峰值/资源链/匹配核验与库存支配论证\\\bottomrule
\end{tabularx}
\end{table}

\textbf{检阅版来源核验状态。}现有记录没有确认上述补充归档的作者、是否属于本队、公开出处及使用许可；归档中未发现明确许可证。因此，本稿披露已知复用事实，不将继承的初始结构写成独立发现，也不据此宣布来源审查通过。正式提交前须补齐作者与取得时间、团队归属或授权、以及本届使用条件的依据。若不能确认可使用条件，应排除受影响的外部候选并从可使用输入重新求解，而不是只删去本段来源说明。该待核验状态不由文字相似度低或方案数值可行自动消除。

\subsection{历史最好解、有限域进展与本次复算的区别}
历史问题二基线和选定方案可以逐项复放，但本次可用归档未保存种子9201形成该候选的全程改进日志及确切轮数，故未补造历史收敛曲线。第\ref{sec:q2}章的20次独立搜索只对应已声明的共同起点和150轮预算。问题三保存的有限网格和任务池状态用于说明实际搜索范围；有限模型的最优、不可行与限时状态不得推广为全空间结论。110份整数状态和125份线性状态也不是235份相互独立的新方案。

本次补写直接列明了下界主问题、有效投影、完整定价和分支覆盖，以及给定资源顺序的连续时间模型；安全载荷曲线由同一物理输入补算，问题三运输时间表和问题四逐组配置由保存的已验证日程直接生成。对新增表格逐项核对箱号、时间和资源向量，不把排版通过当作重新完成全部历史求解或第三方来源授权。
